\section{Matriz conceptual de protocolos}

No existe un único score que responda todas las preguntas del proyecto. La
\cref{tab:validation-matrix} resume los protocolos principales.

\begin{longtable}{p{0.22\linewidth}p{0.18\linewidth}p{0.22\linewidth}p{0.27\linewidth}}
\caption{Protocolos de validación y su interpretación.}
\label{tab:validation-matrix}\\
\toprule
Protocolo & Tamaño & X objetivo visible & Pregunta principal \\
\midrule
\endfirsthead
\toprule
Protocolo & Tamaño & X objetivo visible & Pregunta principal \\
\midrule
\endhead
state-stratified & 5 folds sobre 138 & no como target fijo en 02--03 & interpolación con estados representados.\\
municipality-grouped & 5 folds sobre 138 & no como target fijo en 02--03 & extrapolación a municipios retenidos.\\
target-matched & 16 $\times$ 41 & sí, bajo contrato transductivo & aproximación primaria a la geometría de las 59 parcelas.\\
state-random & 8 $\times$ 41 & sí & control aleatorio respetando estado.\\
LOSO de pseudo-splits & 16 holdouts & sí & controlar selección de método/peso contra mismo split.\\
\bottomrule
\end{longtable}

\section{Nested CV de Checkpoint 03}

La nested CV separa dos capas:
\[
\text{inner CV}\rightarrow \text{selección de hiperparámetro},
\]
\[
\text{outer holdout}\rightarrow \text{estimación de desempeño}.
\]
Formalmente, para outer fold $f$:
\[
\hat\lambda_f=
\arg\min_{\lambda\in\Lambda}
\widehat R_{\mathrm{inner},f}(\lambda),
\]
\[
\widehat R_{\mathrm{outer},f}
=
R\left(
y_{V_f},
f_{\hat\lambda_f,T_f}(X_{V_f})
\right).
\]
La primera operación no puede observar $y_{V_f}$. Esto incluye selección de expresiones
supervisadas, número de componentes PCA/PLS cuando se tunearon y parámetros de modelos.

\section{Pseudo-competition de Checkpoint 04}

En una pseudo-competición $s$, definimos:
\[
\Train_s=\Train\setminus\Target_s^\star,
\qquad
|\Target_s^\star|=41.
\]
El modelo puede usar
\[
X_{\Train},X_{\Target}
\]
para transformaciones $X$-only globales, y puede usar
\[
X_{\Target_s^\star}
\]
durante el diseño de topología. Sin embargo, el operador supervisado satisface:
\[
\hat f_s
=
A(X_{\Train_s},y_{\Train_s};X_{\Target_s^\star}),
\]
con
\[
y_{\Target_s^\star}
\]
reservado hasta evaluación.

Esta distinción es la razón por la que un PCA transductivo puede ser correcto mientras una
selección de variables por correlación con todas las 138 etiquetas no lo sería.

\section{Matching de pseudo-targets}

El target matching utiliza únicamente descriptores de $X$ y metadata. Sea $q(S)$ una función
de calidad de una máscara candidata $S$ que combina distancia entre perfiles, composición
municipal y penalización de diversidad. La máscara seleccionada es aproximadamente:
\[
S^\star
=
\arg\min_{S:|S|=41}
q(S;X_{\Target}),
\]
entre un conjunto finito de candidatos generados con seed congelada.

La validez del procedimiento depende de no introducir ninguna estadística derivada de
$y_S$ en $q$. Las diferencias de yield de vecinos observadas en 04A pueden usarse para
diagnóstico científico, pero no para escoger qué etiquetas se ocultan.

\section{Leave-one-pseudo-split-out}

Para candidate set $\mathcal M$, la selección de método en holdout $h$ es:
\[
m_h^\star
=
\arg\min_{m\in\mathcal M}
\frac{1}{S-1}
\sum_{s\neq h}\RMSE_{s,m}.
\]
Entonces:
\[
R_h
=
\RMSE_{h,m_h^\star}.
\]
El resumen LOSO es
\[
\bar R_{\mathrm{LOSO}}=\frac1S\sum_h R_h.
\]

Este procedimiento controla el error más directo de escoger el mejor método usando el mismo
split que se puntúa. No elimina toda dependencia porque las máscaras reutilizan parcelas. Su
interpretación correcta es \emph{split-excluded selection}, no nested CV con $S$ datasets
independientes.

\section{Candidate universe}

Una condición sutil es que $\mathcal M$ debe ser el mismo universo de métodos que la regla
final pretende utilizar. Si el LOSO puede escoger entre 626 métodos pero la regla real sólo
puede escoger entre diez finalistas, la aparente performance del selector y la regla aplicada
no son el mismo procedimiento.

Este caveat motivó no promover el routing LOSO de 04D.1 pese a su RMSE medio 0.4858. En 04C.1
y 04F el candidate universe se redujo explícitamente antes de seleccionar:
\begin{itemize}
\item 04C.1: Local, CatBoost estable y tres blends Local/CatBoost;
\item 04F: ocho pesos Local/Graph predeclarados.
\end{itemize}

\section{Lectura de diferencias pequeñas}

Para dos métodos $a,b$:
\[
\Delta=\overline{\RMSE}_a-\overline{\RMSE}_b.
\]
Cuando $|\Delta|$ es del orden de $10^{-3}$ o menor, el proyecto exige evidencia adicional:
\begin{enumerate}
\item estabilidad entre splits;
\item pooled RMSE compatible;
\item selección LOSO que no revierta la ventaja;
\item ausencia de un incremento desproporcionado de complejidad;
\item comportamiento razonable en stress protocols.
\end{enumerate}

Este criterio evita convertir ruido de selección en una narrativa de mejora. El caso extremo es
Local+CatBoost10: $\Delta$ en mean RMSE es de apenas $3.5\times10^{-6}$ y no se sostiene en
pooled ni LOSO.

\section{Protocolos de estrés versus objetivo primario}

Municipality-grouped no se optimiza como si fuera el score oculto real; se usa para detectar
fragilidad. Una regla puede ser aceptable si:
\[
R_{\mathrm{target}} \text{ es muy bajo}
\]
y
\[
R_{\mathrm{grouped}} \text{ no introduce un riesgo incompatible},
\]
aun si no es el mínimo grouped.

Graph04D se conserva precisamente por esta razón: no gana el objetivo primario, pero ofrece
una referencia robusta que permite detectar parcelas donde Local puede ser especialmente
sensible a la disponibilidad de vecinos.
